% Original: PDF strana 27, gornji slajd.
\slajdid{02-s053}
\begin{frame}[label=02-s053]{Povezivanje površina: uklanjanje lažne nule}
\setlength{\parskip}{0pt}
\begin{columns}[T,onlytextwidth]
\begin{column}{.31\textwidth}\centering
\Oznaka{Bez zajedničkog polja}\vspace{1mm}
% Izvorne karte su očuvane; formula je usklađena po korisnikovoj odluci.
\begingroup
\fontsize{16.364}{18.2}\selectfont
\renewcommand{\small}{\fontsize{16.364}{18.2}\selectfont}
\scalebox{.55}{%
\begin{karnaugh-map}(label=middle,implicantcolors={DEcrvena,DEcrvena,DEcrvena})[4][2][1][$A$][$B$][$C$]
\minterms{0,1,3,7}\maxterms{2,4,5,6}
{\tikzset{every path/.append style={dashed}}\implicant{0}{1}}
{\tikzset{every path/.append style={dashed}}\implicant{3}{7}}
\end{karnaugh-map}}
\endgroup

\end{column}
\begin{column}{.63\textwidth}
% element: 02-s053:e01
Dve minimizovane površine imaju zajedničku stranu, ali \alert{\textbf{nemaju zajednička polja}}: na prvoj je \(B=0\), na drugoj \(B=1\).\par\vspace{1mm}
Susedni prelaz \(B:1\to0\), pri \(C=0\), \(A=1\), menja samo jednu promenljivu. Različita kašnjenja mogu izazvati \alert{\textbf{lažnu nulu}}, kao na vremenskom dijagramu.
\end{column}
\end{columns}
\par\vspace{1mm}
\begin{columns}[T,onlytextwidth]
\begin{column}{.31\textwidth}\centering
\Oznaka{Povezujuća površina}\vspace{1mm}
% element: 02-s053:e02
% Izvorne karte su očuvane; formula je usklađena po korisnikovoj odluci.
\begingroup
\fontsize{16.364}{18.2}\selectfont
\renewcommand{\small}{\fontsize{16.364}{18.2}\selectfont}
\scalebox{.55}{%
\begin{karnaugh-map}(label=middle,implicantcolors={DEcrvena,DEcrvena,DEcrvena})[4][2][1][$A$][$B$][$C$]
\minterms{0,1,3,7}\maxterms{2,4,5,6}
{\tikzset{every path/.append style={dashed}}\implicant{0}{1}}
{\tikzset{every path/.append style={dashed}}\implicant{3}{7}}
{\tikzset{every path/.append style={solid,line width=1.2pt}}\implicant{1}{3}}
\end{karnaugh-map}}
\endgroup

\end{column}
\begin{column}{.63\textwidth}
{\centering\(F=\bar C\bar B+BA+\textcolor{DEcrvena}{\bar C A}\)\par}\vspace{1mm}
\alert{\textbf{Odstupamo od minimalne forme}}: dodatna površina ima zajednička polja sa obe početne.\par\vspace{1mm}
Član \(\bar C A\) sprečava lažnu nulu pri ovom prelazu, bez promene stabilnih vrednosti funkcije.
\end{column}
\end{columns}
\end{frame}
